% Ley de Fourier
\begin{equation} \label{eq:01}
    q = -kA \frac{dT}{dx}
\end{equation}

\begin{equation} \label{eq:02}
    \rho c_p \frac{\partial T}{\partial t} = \nabla \cdot (k \nabla T) + \dot{q} 
\end{equation}

\begin{equation}
    P = \sigma T^{4}
\end{equation}

%Ley de Planck
\begin{equation}
    E_{\lambda}(\lambda, T) = 
\frac{2\pi h c^{2}}{\lambda^{5}} 
\frac{1}{e^{\frac{hc}{\lambda k T}} - 1}
\end{equation}

%Ley desplazamiento de Wien
\begin{equation}
    \lambda_{\text{max}} T = b
\end{equation}

\begin{equation}
    \alpha + \rho + \tau = 1
\end{equation}

\begin{equation}
    \varepsilon = \alpha
\end{equation}

\begin{equation}
    A_i F_{ij} = A_j F_{ji}
\end{equation}

\begin{equation}
    \sum_{j} F_{ij} = 1
\end{equation}

\begin{equation}
    C_i \frac{dT_i}{dt}
= 
\sum_{j \neq i} K_{ij} (T_j - T_i)
+
\sum_{j \neq i} R_{ij} (T_j^4 - T_i^4)
+
Q_i
\end{equation}

% Son las ecuaciones que aparecen en el marco teorico del TFG de carolina